\SourceLeaf{73}
It follows directly that

\[
S - \frac{dT}{dx} = \left(\frac{\partial W}{\partial r}\right),
\]
\[
R - \frac{dS}{dx} + \frac{d^2T}{dx^2} = \left(\frac{\partial W}{\partial q}\right),
\]
\[
Q - \frac{dR}{dx} + \frac{d^2S}{dx^2} - \frac{d^3T}{dx^3} = \left(\frac{\partial W}{\partial p}\right),
\]
\[
P - \frac{dQ}{dx} + \frac{d^2R}{dx^2} - \frac{d^3S}{dx^3} + \frac{d^4T}{dx^4} = \left(\frac{\partial W}{\partial y}\right),
\]
\[
N - \frac{dP}{dx} + \frac{d^2Q}{dx^2} - \text{etc.} = 0.
\]

The converse theorem can easily be proved from the same principles.

\noindent (Sept.\ 1800.)

\par\smallskip\hrulefill\smallskip\par

\Needspace{10\baselineskip}\noindent\hspace*{\parindent}The theorem states that the condition

\[
\left[N - \frac{dP}{dx} + \frac{d^2Q}{dx^2} - \frac{d^3R}{dx^3} + .\,.\,. = 0\right]
\]

\noindent holds when $Vdx$ is an exact differential. I gave the proof in my copy of \textsc{Euler}’s integral calculus [Volume III, page 518]. The proof of the [converse] theorem rests on the following points.

Since the expression $N - \frac{dP}{dx} + \frac{ddQ}{dx^2} - \text{etc.}$ is derived in a definite way from the nature of the function $V$, we denote it by $fV$, and similarly denote $P - \frac{dQ}{dx} + \frac{ddR}{dx^2} - \text{etc.}$ by $f'V$, and so on.

Set $\int f'V  .  dy = W^0$ by integrating with only $y$ treated as variable, and set the exact differential $dW^0 = V^0dx$. Then\HistoricalNote{*)}{*)}{Both conclusions are connected with the proof at the cited place in \textsc{Euler}’s integral calculus. With $W^0$ in place of $W$, it gives $f'V^0 = \dfrac{\partial W^0}{\partial y}$, hence $= f'V$, and $fV^0 = 0$, hence $= fV$.}

\[
\left.\begin{array}{l}
f'V^0 = f'V \\
fV^0 = fV
\end{array}\right\}
\text{ and hence }
\begin{array}{l}
f'(V - V^0) = 0, \\
f(V - V^0) = 0.
\end{array}
\]

\noindent Now since

\[
\frac{\partial(V - V^0)}{\partial y} = f(V - V^0) + \frac{df'(V - V^0)}{dx}
\]

\noindent [holds,. we also have $\dfrac{\partial(V - V^0)}{\partial y} = 0$] and $V - V^0$ contains no $y$.

